\documentclass[11pt]{article}
\usepackage{amsmath,amsthm,amssymb}
\usepackage[margin=1in]{geometry}
\usepackage{hyperref}
\usepackage{booktabs}
\usepackage{enumitem}

\newtheorem{theorem}{Theorem}
\newtheorem{lemma}[theorem]{Lemma}
\newtheorem{proposition}[theorem]{Proposition}
\newtheorem{corollary}[theorem]{Corollary}
\newtheorem{conjecture}[theorem]{Conjecture}
\theoremstyle{definition}
\newtheorem{definition}[theorem]{Definition}
\newtheorem{remark}[theorem]{Remark}

\newcommand{\Hq}{H_q}
\newcommand{\Sn}{S_n}
\newcommand{\Vlam}{V^\lambda}
\newcommand{\Om}{\Omega}
\newcommand{\Omtail}{\Omega^{(\lambda)}_{\mathrm{tail}}}
\newcommand{\hatl}{\widehat\lambda}
\newcommand{\elh}{\ell(\hatl)}
\newcommand{\Rp}{R'}
\newcommand{\SYT}{\mathrm{SYT}}
\newcommand{\Eplus}[1]{E^+_{#1}}
\newcommand{\Eminus}[1]{E^-_{#1}}
\newcommand{\Vnu}{V^\nu}
\newcommand{\Qsym}{\mathbb{Q}(q)}

\title{Strong Pillar~1 at four tail positions of $(2,2,1,1,1)$:\\
       operator vanishing via a row-1-pair embedding}
\author{Clio}
\date{28 May 2026 (third prove session)}

\begin{document}
\maketitle

\begin{abstract}
Pursuing the strong per-subset operator vanishing conjecture for $\lambda = (2,2,1^m)$,
we prove that the factor-removed tail chains $\Omtail^{[\widehat{\{p\}}]}$ act as the
\emph{zero operator} on $V^\lambda$ for the four positions
$p \in \{12,13,16,17\}$ at $\lambda = (2,2,1,1,1)$. This complements the two structurally
proved positions ($p=10,15$) from~\cite{persubset}, leaving five open positions
($p \in \{11,14,18,19,20\}$) where Strong Pillar~1 ($\subseteq \Eminus 1$) is verified
computationally but the structural proof remains incomplete. The new identities are:

\begin{enumerate}[label=(K\arabic*),topsep=0.2em,itemsep=0.1em]
\item $S_1 S_2 = q \cdot I$ on $\Eplus 1|_{\Vlam}$, by direct Hoefsmit computation
   using the unique axial distance $d=-2$.
\item $S_1 \Rp_n = q \cdot S_{n-1} S_{n-2} \cdots S_3 \cdot S_1$ as operators on
   $\Vlam$, for any $n \ge 4$, by commuting $S_1$ to the right and applying (K1).
\item Lemma~A: $S_1 \Rp_7 \cdot \Vlam \subseteq \Eminus 3 \cap \Eminus 4|_{\Vlam}$,
   by Pillar~1 applied to the embedded sub-representation
   $\Eplus 1|_{\Vlam} \cong V^{(2,1,1,1)}$.
\item Lemma~B: $(\Rp_5)^2 \cdot \Vlam \subseteq \Eminus 6|_{\Vlam}$, by $S_6$-branching
   to the $V^{(2,2,1)}$ summand where the entry pair $\{6,7\}$ is forced into column~1.
\item Lemma~C: $(\Rp_4)^2 \cdot \Vlam \subseteq \Eminus 5|_{\Vlam}$, by the analogous
   $S_5$-branching to the $V^{(2,2)}$-summand where the entry pair $\{5,6\}$ is forced
   into column~1.
\end{enumerate}

Positions $12$ and $13$ vanish via Lemma~A; positions $16$ and $17$ vanish via
Lemmas~B and~C respectively. Computational verification confirms operator vanishing
for all four positions at exact rationals (representative value $q=5/7$), and the
empirical fact $S_1 \Omtail^{[\widehat{\{p\}}]} = 0$ on $\Vlam$ holds for all $11$
tail positions, suggesting a uniform Strong Pillar~1 statement that remains to be
proved at the five open positions.
\end{abstract}

\section{Setup and statement}

We use the conventions of \cite{strong, persubset} throughout. $\Hq(\Sn)$ is the
Iwahori--Hecke algebra over $\Qsym$ with generators $T_1,\ldots,T_{n-1}$ satisfying
$T_i^2 = (q-1)T_i + q$ and the braid/commutation relations. Set $S_i := T_i + 1$;
then $S_i^2 = (q+1)S_i$, so $S_i$ has eigenvalues $0, q+1$, and we write
$\Eplus i = \ker(T_i - q) = \mathrm{Im}(S_i)$, $\Eminus i = \ker(S_i)$.

For $k \ge 1$, $\Rp_k := S_{k-1} S_{k-2} \cdots S_1$ (with $\Rp_1 := I$).
The staircase chain for $\Sn$ is $\Om := \Rp_2 \Rp_3 \cdots \Rp_n$, the reduced word
for the longest element of $\Sn$, of total length $N = \binom{n}{2}$.

For $\lambda = (\hatl, 1^q) \vdash n$ in the Pillar~1 regime
($\elh \ge 2$, every $\hatl_r \ge 2$, $q \ge |\hatl| - 2\elh + 2$), set
$\ell := \elh + q$ and $\tau := q + 2\elh - 1 - |\hatl|$, and split
$\Om = R \cdot \Omtail$ with
\[
   R = \Rp_2 \cdots \Rp_\ell, \qquad \Omtail = \Rp_{\ell+1} \cdots \Rp_n.
\]
The Pillar~1 meta-theorem~\cite{pillar1} states
$\Omtail \cdot \Vlam \subseteq \bigcap_{j=1}^\tau \Eminus j|_{\Vlam}$.

\paragraph{Per-subset operator and Strong Pillar~1.}
For each tail position $p \in \{\binom{\ell}{2}, \ldots, N-1\}$ corresponding to a
staircase letter $i_p$, define the \emph{factor-removed tail chain}
\[
   \Omtail^{[\widehat{\{p\}}]} \;:=\; \prod_{k \in \mathrm{tail},\, k \ne p} S_{i_k},
\]
preserving the staircase order.
The \emph{Strong Pillar~1 property at level~$1$}~\cite{persubset}
is the statement
\begin{equation}\label{eq:strong-pillar1-1}
   \Omtail^{[\widehat{\{p\}}]} \cdot \Vlam \;\subseteq\;
   \Eminus 1|_{\Vlam} \qquad \forall\, p \in \mathrm{tail}.
\end{equation}
This is the key uniform statement we ultimately want. As shown
in~\cite[Thm.~\ref{thm:per-subset-from-strong}]{persubset}, \eqref{eq:strong-pillar1-1}
implies that the per-subset operator $M_T(\{p\})$ satisfies
$M_T(\{p\}) \Vlam \subseteq \Eminus 1$, which combined with prefix kill gives
$M(S) = 0$ for $|S| = 1$.

\paragraph{Target.}
We focus on $\lambda = (2,2,1,1,1)$, so $\hatl = (2,2)$, $\elh = 2$, $q = 3$,
$\ell = 5$, $\tau = 2$, $n = 7$, $N = \binom{7}{2} = 21$. The tail is positions
$10$--$20$ with staircase letters $(5,4,3,2,1; 6,5,4,3,2,1)$.

\begin{theorem}[Operator vanishing at four positions]\label{thm:main}
For $\lambda = (2,2,1,1,1)$, the factor-removed tail chains
$\Omtail^{[\widehat{\{p\}}]}$ act as the zero operator on $\Vlam$ for
$p \in \{12, 13, 16, 17\}$.
\end{theorem}

\noindent
This is strictly stronger than Strong Pillar~1 \eqref{eq:strong-pillar1-1} at
these four positions (which would only require containment in $\Eminus 1$).
The proof proceeds via three structural lemmas (A, B, C), each reducing to
Pillar~1 on a sub-shape via either embedding or branching.

\section{Key identities}

\begin{lemma}[K1: row-1-pair absorption]\label{lem:K1}
On $\Vlam$ for $\lambda = (2,2,1,1,1)$,
\[
   S_1 S_2 \;=\; q \cdot I \qquad \text{on}\; \Eplus 1|_{\Vlam}.
\]
Equivalently, $(S_1 S_2 - q I) \cdot S_1 = 0$ as operators on $\Vlam$.
\end{lemma}

\begin{proof}
$\Eplus 1|_{\Vlam} = \ker(T_1 - q)|_{\Vlam} = \mathrm{Im}(S_1|_{\Vlam})$, of dimension~$4$,
spanned by Hoefsmit basis vectors $e_T$ for SYTs $T$ with $1, 2$ in the same row.
Since $\lambda_1 = 2$ and the cell $(2,2)$ is below $(1,2)$, this forces $1$ at
$(1,1)$ and $2$ at $(1,2)$, with entry $3$ at the only legal position $(2,1)$.

For such $T$, the axial distance of $T$ at $i=2$ is
$d_T(2) = c_T(3) - c_T(2) = (1 - 2) - (2 - 1) = -2$, never $\pm 1$. The Hoefsmit
$2 \times 2$ block at position~$2$ thus reads
\[
   T_2 e_T \;=\; A_{-2} e_T + e_{s_2 T},
   \qquad A_{-2} = \frac{q-1}{1 - q^{-(-2)}} = \frac{q-1}{1-q^2} = -\frac{1}{q+1},
\]
and the swapped tableau $s_2 T$ has $1$ at $(1,1)$, $2$ at $(2,1)$, so
$s_2 T \in \Eminus 1|_{\Vlam}$ (axial dist $d_{s_2 T}(1) = -1$, same-column).
Hence
\[
   S_2 e_T \;=\; \frac{q}{q+1}\, e_T + e_{s_2 T},
\]
with the first summand in $\Eplus 1$ and the second in $\Eminus 1$. Apply $S_1$:
$S_1 e_T = (q+1)\,e_T$ since $e_T \in \Eplus 1$, and $S_1 e_{s_2 T} = 0$ since
$s_2 T \in \Eminus 1$. Therefore
\[
   S_1 S_2 e_T \;=\; \frac{q}{q+1} (q+1)\, e_T + 0 \;=\; q\, e_T.
\]
By linearity, $S_1 S_2 = q I$ on $\Eplus 1|_{\Vlam}$. Since
$\Eplus 1|_{\Vlam} = \mathrm{Im}(S_1|_{\Vlam})$, multiplying on the right by $S_1$
yields the operator identity $(S_1 S_2 - qI) S_1 = 0$.\qed
\end{proof}

\begin{lemma}[K2: $S_1 \Rp_n$ collapse]\label{lem:K2}
For any $n \ge 4$, $S_1 \Rp_n = q \cdot S_{n-1} S_{n-2} \cdots S_3 \cdot S_1$ as
operators on $\Vlam$ ($\lambda = (2,2,1,1,1)$).
\end{lemma}

\begin{proof}
$\Rp_n = S_{n-1} S_{n-2} \cdots S_2 S_1$. Since $S_1$ commutes with $S_j$ for $j \ge 3$
($|1-j| \ge 2$), we can pull $S_1$ rightward past $S_{n-1}, \ldots, S_3$:
\[
   S_1 \Rp_n \;=\; S_{n-1} S_{n-2} \cdots S_3 \cdot S_1 S_2 S_1.
\]
For any $v \in \Vlam$, $S_1 v \in \mathrm{Im}(S_1|_{\Vlam}) = \Eplus 1|_{\Vlam}$.
By Lemma~\ref{lem:K1} applied to $w = S_1 v$, $S_1 S_2 w = q w$, i.e.\
\[
   S_1 S_2 S_1 v \;=\; S_1 S_2 (S_1 v) \;=\; q (S_1 v) \;=\; q S_1 v.
\]
Therefore $S_1 S_2 S_1 = q S_1$ as operators on $\Vlam$, and
\[
   S_1 \Rp_n \;=\; S_{n-1} \cdots S_3 \cdot S_1 S_2 S_1
   \;=\; q \cdot S_{n-1} \cdots S_3 \cdot S_1.\qed
\]
\end{proof}

\section{Structural lemmas}

\begin{lemma}[Lemma~A]\label{lem:A}
$S_1 \Rp_7 \cdot \Vlam \subseteq \Eminus 3 \cap \Eminus 4|_{\Vlam}$.
\end{lemma}

\begin{proof}
By Lemma~\ref{lem:K2} (with $n = 7$),
\[
   S_1 \Rp_7 \cdot v \;=\; q \cdot S_6 S_5 S_4 S_3 \cdot S_1 v
\]
for every $v \in \Vlam$. So
\[
   S_1 \Rp_7 \cdot \Vlam \;=\; q \cdot S_6 S_5 S_4 S_3 \cdot \Eplus 1|_{\Vlam}.
\]

The subspace $\Eplus 1|_{\Vlam}$ is spanned by Hoefsmit vectors $e_T$ with $1, 2$
at $(1,1), (1,2)$; the remaining entries $3,\ldots,7$ then fill the sub-shape
$\nu := (2,1,1,1)$ (cells $(2,1), (2,2), (3,1), (4,1), (5,1)$ of $\lambda$),
according to a standard tableau $\widetilde T$ of $\nu$ with entries $\{3,\ldots,7\}$.
Equivalently, with the relabeling $i \leftrightarrow i - 2$ for $i \ge 3$ in $T$,
the map $T \mapsto \widetilde T$ gives a bijection between SYT of $\lambda$ in
$\Eplus 1|_{\Vlam}$ and SYT of $\nu = (2,1,1,1)$.

For $i \ge 3$, the Hecke generator $T_i$ on $e_T$ depends only on the positions of
$i$ and $i+1$ in $T$, both of which lie in $\nu$. Therefore the action of
$T_3, T_4, T_5, T_6$ on $\Eplus 1|_{\Vlam}$ is intertwined by this relabeling
bijection with the action of $\widetilde T_1, \widetilde T_2, \widetilde T_3,
\widetilde T_4$ on $V^\nu$ in the standard Hoefsmit basis.

Under this identification, the operator $S_6 S_5 S_4 S_3$ on $\Eplus 1|_{\Vlam}$
becomes $\widetilde S_4 \widetilde S_3 \widetilde S_2 \widetilde S_1 = \widetilde \Rp_5$
on $V^\nu$. The shape $\nu = (2,1^3)$ is a hook with arm $k_h = 2$ and leg
$m_h = 3$. Its Pillar~1 parameters are
$\hatl(\nu) = (2)$, $\elh(\nu) = 1$, $q'(\nu) = 3$, so
$\tau(\nu) = q' + 2 - 1 - 2 = 2$ and $\ell(\nu) = 4$, with $\Omtail^{(\nu)} = \widetilde \Rp_5$
($\ell + 1 = 5 = n$, a single block).
Pillar~1 (Theorem~\ref{thm:pillar1}) applied to $\nu$ gives
\[
   \widetilde \Rp_5 \cdot V^\nu \;\subseteq\; \widetilde{\Eminus 1} \cap \widetilde{\Eminus 2}|_{V^\nu}.
\]
By the relabeling, $\widetilde{\Eminus 1} = \Eminus 3$ and $\widetilde{\Eminus 2} = \Eminus 4$
inside $\Vlam$. Therefore
\[
   S_6 S_5 S_4 S_3 \cdot \Eplus 1|_{\Vlam}
   \;\subseteq\; \Eminus 3 \cap \Eminus 4|_{\Vlam},
\]
and $S_1 \Rp_7 \cdot \Vlam$ is a scalar multiple of this, hence contained in the
same intersection.\qed
\end{proof}

\begin{lemma}[Lemma~B]\label{lem:B}
$(\Rp_5)^2 \cdot \Vlam \subseteq \Eminus 6|_{\Vlam}$.
\end{lemma}

\begin{proof}
$\Rp_5 \in \Hq(S_5)$, so $(\Rp_5)^2 \in \Hq(S_5) \subseteq \Hq(S_6)$ preserves the
$\Hq(S_6)$-branching $\Vlam \downarrow S_6 = W_1 \oplus W_2$, where
$W_1 = V^{(2,1,1,1,1)}$ (the summand with $7$ at $(2,2)$ in $\Vlam$) and
$W_2 = V^{(2,2,1,1)}$ ($7$ at $(5,1)$).

\paragraph{On $W_1$:}
$W_1 \downarrow S_5 = V^{(1^5)} \oplus V^{(2,1,1,1)}$ corresponding to removable
cells $(1,2)$ and $(5,1)$ of $(2,1,1,1,1)$.
\begin{itemize}[topsep=0.2em,itemsep=0.1em]
\item On the sign-representation summand $V^{(1^5)}$: every $S_i$ acts as $0$, so
   $\Rp_5 \equiv 0$ there, hence $(\Rp_5)^2 V^{(1^5)} = 0$.
\item On $V^{(2,1,1,1)}$: this is a hook (arm $2$, leg $3$) viewed as an $\Hq(S_5)$-module.
   By Pillar~1 with $\hatl = (2)$, $q' = 3$, $\tau = 2$, the chain $\Rp_5$ (= the
   Pillar~1 tail $\Omtail^{(2,1,1,1)}$ in $S_5$) sends $V^{(2,1,1,1)}$ into
   $\Eminus 1|_{V^{(2,1,1,1)}}$. Then $\Rp_5$ applied again ends in $S_1$, which
   annihilates $\Eminus 1$. Hence $(\Rp_5)^2 V^{(2,1,1,1)} = 0$.
\end{itemize}
Combining, $(\Rp_5)^2 W_1 = 0$.

\paragraph{On $W_2$:}
$W_2 \downarrow S_5 = V^{(2,1,1,1)} \oplus V^{(2,2,1)}$ corresponding to removable
cells $(2,2)$ and $(4,1)$ of $(2,2,1,1)$.
\begin{itemize}[topsep=0.2em,itemsep=0.1em]
\item On $V^{(2,1,1,1)}$: same argument as above, $(\Rp_5)^2 V^{(2,1,1,1)} = 0$.
\item On $V^{(2,2,1)}$: in this $\Hq(S_5)$-summand, the entries $6, 7$ are
   positioned at $(4,1)$ and $(5,1)$ respectively (from the iterated branching trace).
   In particular, $6$ and $7$ lie in the \emph{same column}, so the Hoefsmit
   action of $T_6$ is $T_6 e_T = -e_T$ (axial distance $-1$), giving
   $S_6 e_T = (T_6 + I) e_T = 0$ on the entire $V^{(2,2,1)}$-summand.
\end{itemize}
Therefore $S_6 \cdot (\Rp_5)^2 W_2 \subseteq S_6 V^{(2,2,1)} = 0$.

Combining $W_1$ and $W_2$: $S_6 \cdot (\Rp_5)^2 \Vlam = 0$, i.e.,
$(\Rp_5)^2 \Vlam \subseteq \ker S_6 = \Eminus 6$.\qed
\end{proof}

\begin{lemma}[Lemma~C]\label{lem:C}
$(\Rp_4)^2 \cdot \Vlam \subseteq \Eminus 5|_{\Vlam}$.
\end{lemma}

\begin{proof}
Analogous to Lemma~\ref{lem:B}. $\Rp_4 \in \Hq(S_4) \subseteq \Hq(S_5)$ preserves the
$\Hq(S_5)$-branching of $\Vlam$. By the iterated branching trace,
\[
   \Vlam \downarrow S_5 \;=\; V^{(1^5)} \oplus 2\, V^{(2,1,1,1)} \oplus V^{(2,2,1)},
\]
where the $V^{(2,2,1)}$-summand has entries $5, 6, 7$ at positions $(3,1), (4,1), (5,1)$
respectively (deepest column-1 stack). On the sign and hook summands, the previous
argument gives $(\Rp_4)^2 = 0$:
\begin{itemize}[topsep=0.2em,itemsep=0.1em]
\item $V^{(1^5)}$: trivially.
\item $V^{(2,1,1,1)}$: $\Rp_4$ is the Pillar~1 tail in $\Hq(S_4)$ acting on
   $V^{(2,1,1,1)} \downarrow S_4$, and using hook column-1 plus $S_1$-end kill,
   $(\Rp_4)^2 = 0$.
\item $V^{(2,2,1)}$: this is an $\Hq(S_5)$-module restricted to $\Hq(S_4)$:
   $V^{(2,2,1)} \downarrow S_4 = V^{(2,1,1)} \oplus V^{(2,2)}$ from removable cells
   $(2,2), (3,1)$. On $V^{(2,1,1)}$ (hook), $(\Rp_4)^2 = 0$ by the same chain
   argument. On $V^{(2,2)}$: in this sub-summand the entries $5, 6, 7$ are at
   positions $(3,1), (4,1), (5,1)$ in $\lambda$; in particular $5, 6$ are
   same-column.
\end{itemize}
On the $V^{(2,2)}$-summand: $T_5 e_T = -e_T$ (same-column for $5, 6$), so
$S_5 e_T = 0$ there. Therefore $S_5 \cdot (\Rp_4)^2 \Vlam = 0$, equivalently
$(\Rp_4)^2 \Vlam \subseteq \Eminus 5|_{\Vlam}$.\qed
\end{proof}

\section{Proof of Theorem~\ref{thm:main}}

\subsection{Positions $12$ and $13$ via Lemma~A}

\begin{lemma}[Position $12$]\label{lem:pos12}
$\Omtail^{[\widehat{\{12\}}]} = 0$ on $\Vlam$.
\end{lemma}

\begin{proof}
Position $12$ is the third letter of $\Rp_6 = S_5 S_4 S_3 S_2 S_1$, which is $S_3$.
Removing it gives $\Rp_6^{[\widehat{S_3}]} = S_5 S_4 S_2 S_1$, and
\[
   \Omtail^{[\widehat{\{12\}}]} \;=\; (S_5 S_4 S_2 S_1) \cdot \Rp_7.
\]
Reading the chain $S_5 S_4 S_2 S_1 \Rp_7$ from the right (operator application
order), the factors apply in turn: $\Rp_7$, $S_1$, $S_2$, $S_4$, $S_5$.

For any $v \in \Vlam$, set $w := \Rp_7 v$ and consider the partial product
$S_2 S_1 \cdot w = S_2 (S_1 w)$. By Lemma~\ref{lem:A}, $S_1 \Rp_7 \Vlam \subseteq
\Eminus 4|_{\Vlam}$, so $S_1 w \in \Eminus 4$. Since $|2 - 4| = 2$, $S_2$ commutes
with $T_4$, hence preserves $\Eminus 4 = \ker S_4$. Therefore
\[
   S_2 S_1 \Rp_7 v \;\in\; \Eminus 4|_{\Vlam}.
\]
Applying $S_4$: $S_4 \cdot S_2 S_1 \Rp_7 v = 0$. The two remaining left factors $S_5$
do not change this:
\[
   \Omtail^{[\widehat{\{12\}}]} v \;=\; S_5 \cdot S_4 \cdot S_2 \cdot S_1 \cdot \Rp_7 v
   \;=\; S_5 \cdot 0 \;=\; 0.\qed
\]
\end{proof}

\begin{lemma}[Position $13$]\label{lem:pos13}
$\Omtail^{[\widehat{\{13\}}]} = 0$ on $\Vlam$.
\end{lemma}

\begin{proof}
Position $13$ is the fourth letter of $\Rp_6$, which is $S_2$. Removing it gives
$\Rp_6^{[\widehat{S_2}]} = S_5 S_4 S_3 S_1$, and
\[
   \Omtail^{[\widehat{\{13\}}]} \;=\; (S_5 S_4 S_3 S_1) \cdot \Rp_7.
\]
For any $v \in \Vlam$, set $w := \Rp_7 v$. By Lemma~\ref{lem:A}, $S_1 w \in \Eminus 3$.
Applying $S_3$: $S_3 \cdot S_1 w = 0$. Therefore
\[
   \Omtail^{[\widehat{\{13\}}]} v \;=\; S_5 S_4 \cdot S_3 \cdot S_1 \cdot \Rp_7 v
   \;=\; S_5 S_4 \cdot 0 \;=\; 0.\qed
\]
\end{proof}

\subsection{Positions $16$ and $17$ via Lemmas~B and~C}

\begin{lemma}[Position $16$]\label{lem:pos16}
$\Omtail^{[\widehat{\{16\}}]} = 0$ on $\Vlam$.
\end{lemma}

\begin{proof}
Position $16$ is the second letter of $\Rp_7 = S_6 S_5 S_4 S_3 S_2 S_1$, which is $S_5$.
Removing it leaves $\Rp_7^{[\widehat{S_5}]} = S_6 S_4 S_3 S_2 S_1 = S_6 \Rp_5$,
where $S_6$ commutes with each factor of $\Rp_5 = S_4 S_3 S_2 S_1$ (all distances
$\ge 2$). Therefore
\[
   \Omtail^{[\widehat{\{16\}}]} \;=\; \Rp_6 \cdot S_6 \Rp_5
   \;=\; \Rp_6 \cdot \Rp_5 \cdot S_6.
\]
Using $\Rp_6 = S_5 \Rp_5$:
\[
   \Rp_6 \Rp_5 \;=\; S_5 \cdot (\Rp_5)^2.
\]
Combining and commuting $S_6$ left past $(\Rp_5)^2$ (which contains only
$S_4, S_3, S_2, S_1$):
\[
   \Omtail^{[\widehat{\{16\}}]} \;=\; S_5 \cdot (\Rp_5)^2 \cdot S_6
   \;=\; S_5 S_6 \cdot (\Rp_5)^2.
\]
For any $v \in \Vlam$, by Lemma~\ref{lem:B}, $(\Rp_5)^2 v \in \Eminus 6|_{\Vlam}
= \ker S_6$. Hence $S_6 \cdot (\Rp_5)^2 v = 0$, and
$\Omtail^{[\widehat{\{16\}}]} v = S_5 \cdot 0 = 0$.\qed
\end{proof}

\begin{lemma}[Position $17$]\label{lem:pos17}
$\Omtail^{[\widehat{\{17\}}]} = 0$ on $\Vlam$.
\end{lemma}

\begin{proof}
Position $17$ is the third letter of $\Rp_7$, namely $S_4$. Removing it leaves
$\Rp_7^{[\widehat{S_4}]} = S_6 S_5 S_3 S_2 S_1 = S_6 S_5 \cdot \Rp_4$
(where $\Rp_4 = S_3 S_2 S_1$). Therefore
\[
   \Omtail^{[\widehat{\{17\}}]} \;=\; \Rp_6 \cdot S_6 S_5 \Rp_4.
\]
Commute $\Rp_4$ left through $S_5$ and $S_6$ (distances $\ge 2$ from $\{3, 2, 1\}$):
$\Rp_4 \cdot S_5 = S_5 \cdot \Rp_4$, $\Rp_4 \cdot S_6 = S_6 \cdot \Rp_4$. So
\[
   \Omtail^{[\widehat{\{17\}}]} \;=\; \Rp_6 \Rp_4 \cdot S_6 S_5.
\]
Using $\Rp_6 = S_5 S_4 \Rp_4$ (since $\Rp_5 = S_4 \Rp_4$ and $\Rp_6 = S_5 \Rp_5$):
\[
   \Rp_6 \Rp_4 \;=\; S_5 S_4 \cdot \Rp_4 \cdot \Rp_4 \;=\; S_5 S_4 \cdot (\Rp_4)^2,
\]
so
\[
   \Omtail^{[\widehat{\{17\}}]} \;=\; S_5 S_4 \cdot (\Rp_4)^2 \cdot S_6 S_5.
\]
Commute $S_6$ and the trailing $S_5$ leftward past $(\Rp_4)^2$ (all commute):
$S_6 \cdot (\Rp_4)^2 = (\Rp_4)^2 \cdot S_6$, $S_5 \cdot (\Rp_4)^2 = (\Rp_4)^2 \cdot S_5$.
Inserting these in turn:
\[
   \Omtail^{[\widehat{\{17\}}]}
   \;=\; S_5 S_4 \cdot S_6 \cdot (\Rp_4)^2 \cdot S_5
   \;=\; S_5 S_4 S_6 S_5 \cdot (\Rp_4)^2.
\]
For any $v \in \Vlam$, the rightmost factor applied is $(\Rp_4)^2$, and by
Lemma~\ref{lem:C}, $(\Rp_4)^2 v \in \Eminus 5|_{\Vlam} = \ker S_5$. The next
applied factor is $S_5$, which annihilates this:
\[
   \Omtail^{[\widehat{\{17\}}]} v
   \;=\; S_5 S_4 S_6 \cdot S_5 \cdot (\Rp_4)^2 v
   \;=\; S_5 S_4 S_6 \cdot 0 \;=\; 0.\qed
\end{proof}

\section{Computational verification and status of open positions}

The following table summarizes the $11$ tail positions of $\lambda = (2,2,1,1,1)$
($\tau = 2$, tail length $11$). Ranks and containments are computed exactly at
$q = 5/7$ in the Hoefsmit basis. (Equivalent results hold at $q = 2/3$.)

\begin{table}[h]
\centering\renewcommand{\arraystretch}{1.2}
\caption{Strong Pillar~1 status across all $11$ tail positions at
$\lambda = (2,2,1,1,1)$. ``Rank'' is $\mathrm{rk}\,\Omtail^{[\widehat{\{p\}}]}|_{\Vlam}$.
Empirically, $S_1 \cdot \Omtail^{[\widehat{\{p\}}]} = 0$ on $\Vlam$ for all $11$
positions, which is the Strong Pillar~1 statement at level~$1$.}
\begin{tabular}{cccccc}
\toprule
$p$ & letter & rank & $\subseteq \Eminus 1$ & operator vanishing & structural proof \\
\midrule
$10$ & $S_5$ & $1$ & yes & no & \cite[Lem.~\ref{lem:pos10}]{persubset} \\
$11$ & $S_4$ & $1$ & yes & no & open \\
$12$ & $S_3$ & $0$ & yes & \textbf{yes} & Lem.~\ref{lem:pos12} (this paper) \\
$13$ & $S_2$ & $0$ & yes & \textbf{yes} & Lem.~\ref{lem:pos13} (this paper) \\
$14$ & $S_1$ & $1$ & yes & no & open \\
$15$ & $S_6$ & $1$ & yes & no & \cite[Lem.~\ref{lem:pos15}]{persubset} \\
$16$ & $S_5$ & $0$ & yes & \textbf{yes} & Lem.~\ref{lem:pos16} (this paper) \\
$17$ & $S_4$ & $0$ & yes & \textbf{yes} & Lem.~\ref{lem:pos17} (this paper) \\
$18$ & $S_3$ & $1$ & yes & no & open \\
$19$ & $S_2$ & $1$ & yes & no & open \\
$20$ & $S_1$ & $1$ & yes & no & open \\
\bottomrule
\end{tabular}
\end{table}

\paragraph{Open positions.}
For positions $11, 14, 18, 19, 20$ (rank~$1$), the operator
$\Omtail^{[\widehat{\{p\}}]}$ does \emph{not} vanish on $\Vlam$ (its image is a
rank-$1$ subspace contained in $\Eminus 1|_{\Vlam}$), but the structural proof
of Strong Pillar~1 \eqref{eq:strong-pillar1-1} at these positions remains open.
Computationally:
\begin{itemize}[topsep=0.2em,itemsep=0.1em]
\item All five rank-$1$ images $\Omtail^{[\widehat{\{p\}}]} \Vlam$ are supported on
   single SYTs (in the Hoefsmit basis), specifically: $p=10 \to T_9$, $p=11 \to T_5$,
   $p=14 \to T_8$, $p=15 \to T_5$, $p=18 \to T_8$, $p=19 \to T_8$, $p=20 \to T_8$.
   In each case, the supporting SYT has $6$ and $7$ same-column (column~$1$), but
   $1$ and $2$ at $(1,1), (1,3)$ or similar — not the row-$1$ pattern that triggers
   our Lemma~A bypass.
\item The same key identity (Lemma~\ref{lem:K2}) immediately yields
   $S_1 \cdot \Omtail^{[\widehat{\{p\}}]} = 0$ \emph{empirically} for all $11$
   positions, but the structural proof is direct only when the chain can be
   factorized to expose an inner $S_1$-collapse identity (as in Lemmas~A--C).
\end{itemize}

\paragraph{Sketch of a uniform conjectural statement.}
Define
$\mathcal{N}_p := \Omtail^{[\widehat{\{p\}}]} \Vlam$. The empirical data suggests
\begin{equation}\label{eq:uniform}
   S_1 \cdot \mathcal{N}_p = 0 \quad \forall p \in \mathrm{tail},
\end{equation}
i.e., Strong Pillar~1 at level~$1$. By the linear-combination identity
\cite[Prop.~\ref{prop:identity}]{persubset},
$M_T(\{p\}) = \frac{1}{q+1}\Omtail^{[\widehat{\{p\}}]} - \frac{1}{(q+1)^{|\mathrm{tail}|}}\Omtail$
modulo a global scaling, so once \eqref{eq:uniform} is proved, per-subset
Pillar~1 follows for $|S_T| = 1$, closing $4944 - $ (boundary outliers) of the
$M(S) = 0$ cases for $|S| \le 1$ at $m = 3$.

\section{Status}

\paragraph{Established (this paper).}
\begin{itemize}[topsep=0.2em,itemsep=0.1em]
\item Lemma~\ref{lem:K1}: $S_1 S_2 = q I$ on $\Eplus 1|_{\Vlam}$ for
   $\lambda = (2,2,1,1,1)$, via direct Hoefsmit calculation using axial distance $-2$.
\item Lemma~\ref{lem:K2}: $S_1 \Rp_n = q S_{n-1} \cdots S_3 \cdot S_1$ on $\Vlam$
   for $n \ge 4$.
\item Lemma~\ref{lem:A}: $S_1 \Rp_7 \Vlam \subseteq \Eminus 3 \cap \Eminus 4$, via
   embedding of $\Eplus 1|_{\Vlam}$ into the sub-shape $V^{(2,1,1,1)}$ and Pillar~1
   for the hook $(2,1^3)$.
\item Lemma~\ref{lem:B}: $(\Rp_5)^2 \Vlam \subseteq \Eminus 6$, via $S_6$-branching
   and same-column kill on $V^{(2,2,1)}$.
\item Lemma~\ref{lem:C}: $(\Rp_4)^2 \Vlam \subseteq \Eminus 5$, via $S_5$-branching
   to $V^{(2,2)}$.
\item Theorem~\ref{thm:main}: $\Omtail^{[\widehat{\{p\}}]} = 0$ on $\Vlam$ for
   $p \in \{12, 13, 16, 17\}$.
\end{itemize}

\paragraph{Open gaps (precisely stated).}
\begin{itemize}[topsep=0.2em,itemsep=0.1em]
\item Strong Pillar~1 at the five rank-$1$ positions $p \in \{11, 14, 18, 19, 20\}$:
   the image is verified to lie in $\Eminus 1$ at exact rationals, and the empirical
   $S_1 \cdot \Omtail^{[\widehat{\{p\}}]} = 0$ identity holds, but the structural
   proof requires a different mechanism (the leftmost-factor / branching reduction
   used for $p \in \{12, 13, 16, 17\}$ does not directly apply).
\item Per-subset Pillar~1 outliers for $|S_T| = 2$: the three subsets
   $\{11, 12\}, \{12, 15\}, \{15, 16\}$ identified
   in~\cite{strong,persubset}, where $\tau - |S_T| = 0$ and refined Pillar~1
   makes no claim. The inductive-kernel mechanism remains case-by-case.
\item Extension to $\lambda = (2, 2, 1^m)$ for $m \ge 4$: computational sweep
   beyond $m = 3$ and inductive argument across $m$ both open.
\end{itemize}

\paragraph{Methodological note.}
The key insight enabling the proofs of positions $12, 13$ is the
identity $S_1 S_2 = qI$ on $\Eplus 1|_{\Vlam}$, which arises directly from the
Hoefsmit axial distance $d = -2$ for the unique non-degenerate $2 \times 2$ block at
position~$2$ inside $\Eplus 1|_{\Vlam}$. This identity propagates through the
$\Rp_n$-chain via Coxeter commutation, ultimately reducing $S_1 \Rp_n V^\lambda$
to a Pillar~1 computation on the sub-shape $V^{(2, 1^{n-4})}$.
This is the first instance, in our program, of a per-subset chain reduction by
explicit Hoefsmit identity rather than by chain factorization or full branching;
it is likely to extend to other off-hook shapes whose first column hosts a
``row-$1$-pair-plus-sub-hook'' decomposition.

The remaining five open positions all involve dropping a letter from $\Rp_6$ or
$\Rp_7$ at a position where the resulting chain factor structure does not admit a
clean $(\Rp_k)^2$ or Lemma~A reduction. A uniform mechanism — perhaps via the
linear-combination identity of \cite[Prop.~\ref{prop:identity}]{persubset} applied
in tandem with the chain operator algebra — is the natural next target.

\begin{thebibliography}{9}
\bibitem{strong}
   Clio, \emph{Strong operator vanishing of the staircase chain on $V^{(2,2,1^m)}$
   (the $|S|=0$ case of strong per-subset vanishing)}, 28 May 2026.
   \texttt{\textasciitilde/projects/proofs/2026-05-28-strong-per-subset-operator-vanishing.tex}.

\bibitem{persubset}
   Clio, \emph{A linear-combination identity for per-subset Pillar~1:
   reducing strong operator vanishing to a single tail statement}, 28 May 2026.
   \texttt{\textasciitilde/projects/proofs/2026-05-28-per-subset-pillar1-reduction.tex}.

\bibitem{pillar1}
   Clio, \emph{Extended sign-kill for general multi-row $\hatl$: the
   meta-theorem at arbitrary $\elh \ge 2$}, 15 May 2026.
   \texttt{\textasciitilde/projects/proofs/2026-05-15-multi-row-sign-kill-meta.tex}.

\bibitem{sharpness}
   Clio, \emph{Sharpness of the multi-row containment at $\tau+1$:
   $B^{(\lambda)} \cap E^-_{\tau+1}|_{\Vlam} = 0$}, 17 May 2026.
   \texttt{\textasciitilde/projects/proofs/2026-05-17-sharpness-tau-plus-1.tex}.
\end{thebibliography}

\end{document}
